\begin{abstract}
Este trabajo presenta una simulación basada en el modelo de Schelling para estudiar escenarios de segregación residencial en una representación simplificada de Montevideo. El modelo incorpora composición socioeconómica del vecindario, restricciones de tolerancia, precios de vivienda, capacidad económica y permanencia mínima. Se desarrollaron una implementación secuencial reproducible y una implementación paralela híbrida con MPI y OpenMP, equivalentes bit a bit en las configuraciones verificadas. En una medición local de una iteración sobre una grilla de $1024\times640$, la mediana secuencial fue 3,672~s y la mejor configuración híbrida, con dos procesos y cuatro threads por proceso, alcanzó 1,091~s, un speedup de 3,366. El resultado confirma que la búsqueda de destinos admite paralelismo útil, aunque la eficiencia decrece por trabajo replicado, comunicación y costos seriales. La evaluación entre nodos de FING queda separada de estos resultados locales para evitar extrapolaciones entre plataformas.
\end{abstract}

\begin{IEEEkeywords}
Schelling, modelos basados en agentes, segregación residencial, MPI, OpenMP, computación de alta performance.
\end{IEEEkeywords}
